\documentclass[11pt]{article}
\usepackage[hyperref]{acl}
\usepackage{times}
\usepackage{latexsym}
\usepackage[T1]{fontenc}
\usepackage[utf8]{inputenc}
\usepackage{microtype}
\usepackage{inconsolata}
\usepackage{booktabs}
\usepackage{amsmath}
\usepackage{amssymb}
\usepackage{graphicx}

\title{Comparing Chunking Strategies for Retrieval-Augmented Question Answering}

\author{
  Assylkhan Geniyat \\
  Boston University \\
  \texttt{kuromiqo@bu.edu} \And
  Pablo Bello \\
  Boston University \\
  \texttt{pbello@bu.edu} \And
  Jaime Bernal \\
  Boston University \\
  \texttt{jaimebm@bu.edu} \AND
  Batyrkhan Baimukhanov \\
  Boston University \\
  \texttt{batyr@bu.edu}
}

\date{}

\begin{document}
\maketitle

\begin{abstract}
We study how document chunking affects retrieval-augmented question answering in a controlled pipeline where the embedding model, retriever, and generator are held fixed. Four chunking families are compared: fixed-size token windows, recursive chunking, sentence-based chunking, and semantic chunking. Using \texttt{Mistral-7B-Instruct-v0.3} served through vLLM on SCC, we find that \texttt{recursive\_256} is the strongest system on both a 60-question SQuAD v2 subset (60.6 F1) and a 30-question HotpotQA subset (42.6 F1). A parametric-only baseline is substantially weaker, confirming retrieval sensitivity. These findings are preliminary: the subsets are small, and strict metrics may undercount correct answers that differ only in surface form. Nevertheless, the results already show that chunking is a meaningful experimental variable whose best setting may depend on question complexity.
\end{abstract}

\section{Introduction}

Retrieval-augmented generation (RAG) answers questions by retrieving external evidence and passing it to a language model \citep{lewis2020rag}. In practice, retrieval operates over chunks rather than full documents, making chunking a central design choice: it determines what evidence stays together, how likely the retriever is to recover it, and whether the generator receives coherent support or fragmented context.

This project asks: \textit{which chunking strategy maximizes end-to-end QA accuracy in a fixed RAG pipeline, and does the answer differ between local factoid questions and multi-hop questions?}

Formally, let a corpus be $D$, a question be $q$, and a chunking function be $c : D \rightarrow \{d_1, \dots, d_n\}$. After chunking, each chunk is embedded and indexed. At inference time, the system retrieves the top-$K$ chunks for $q$, concatenates them into a prompt, and produces an answer $\hat{a}$. The optimization target is end-to-end QA quality:
\[
c^* = \arg\max_c \; \mathbb{E}_{(q,a)} \left[ F_1(\hat{a}, a) \right].
\]

This problem is interesting because retrieval quality and answer quality do not always move together: a system may retrieve locally relevant text yet still fail if the evidence is incomplete or badly fragmented. We therefore evaluate chunking with both end-to-end QA metrics and retrieval diagnostics.

\section{Related Work}

RAG combines neural retrieval with text generation by conditioning a generator on retrieved passages \citep{lewis2020rag}. Dense Passage Retrieval \citep{karpukhin2020dpr} established dense retrieval as a strong QA backbone and serves as our fixed retriever. We evaluate on SQuAD~2.0 \citep{rajpurkar2018squad2} for local extraction and HotpotQA \citep{yang2018hotpotqa} for multi-hop reasoning.

Recent work motivates our chunking focus. \citet{smith2024chunking} show that retrieval quality varies substantially across chunking policies, NVIDIA argues that chunking should be judged end-to-end \citep{nvidia2024chunking}, and \citet{qu2024semanticchunking} question whether semantic chunking always justifies its cost. We isolate chunking as the only variable inside a reproducible QA pipeline.

\section{Methods}

For each chunking strategy, the pipeline chunks each document, embeds every chunk with the same encoder, builds a FAISS dense index, retrieves the top~4 chunks for each question, and passes them to the same generator. The chunking function is the only variable.

\subsection{Chunking Strategies}

We implement four families:
\begin{itemize}
  \item \textbf{Fixed token chunking:} \texttt{fixed\_128}, \texttt{fixed\_256}, \texttt{fixed\_512}, with $\sim$15\% overlap.
  \item \textbf{Recursive chunking:} \texttt{recursive\_256}, using LangChain's recursive splitter.
  \item \textbf{Sentence chunking:} \texttt{sentence\_256}, using spaCy segmentation with greedy packing to a token budget.
  \item \textbf{Semantic chunking:} \texttt{semantic\_256}, splitting when inter-sentence cosine similarity drops below 0.72, with a minimum-length guard. Chunk length varies with local semantic continuity.
\end{itemize}

We also implemented auxiliary \texttt{chonkie\_recursive\_256} and \texttt{chonkie\_semantic\_256} variants to check whether results are robust across different library implementations of the same chunking idea.

\subsection{Models and Configuration}

The pipeline uses \texttt{all-MiniLM-L6-v2} for embedding, FAISS inner-product search over normalized embeddings for retrieval, and \texttt{Mistral-7B-Instruct-v0.3} served through vLLM on SCC as the generator. Prompting consists of a short extractive QA instruction plus a lightweight answer-compression step for verbose outputs. Generation is limited to 1024 input tokens and 48 new tokens. The broader codebase also supports BM25, hybrid retrieval, and reranking, but the midway analysis keeps the ablation focused on chunking. All hyperparameters are fixed in JSON configs and reused across runs.

\subsection{Baseline}

Our baseline is \textbf{parametric-only QA}: the generator answers without retrieved context. This tests whether retrieval is necessary and provides a lower bound.

\section{Experiments}

\subsection{Datasets}

\begin{itemize}
  \item \textbf{SQuAD v2:} 60 answerable validation questions sampled from a 500-example pool. Paragraphs sharing a title are concatenated into 35 article-level pseudo-documents to make chunking non-trivial.
  \item \textbf{HotpotQA distractor:} 30 validation questions with 300 documents built from distractor and supporting titles.
\end{itemize}

\subsection{Evaluation}

We report Exact Match (EM), token-level F1, Recall@4, Precision@4, average chunk size, and chunk count. For SQuAD, Recall@4 reduces to single-document recall. For HotpotQA, it is the fraction of gold supporting documents recovered in the top~4 chunks, which is more informative than an ``any-hit'' definition for multi-hop questions.

\section{Results}

% Auto-generated by scripts/export_report_tables.py
\begin{table*}[t]
\centering
\small
\begin{tabular}{lrrrrrr}
\toprule
System & EM & F1 & Recall@4 & Precision@4 & Avg chunk tokens & \# chunks \\
\midrule
\texttt{parametric\_only} & 13.3 & 29.2 & 0.0 & 0.0 & -- & -- \\
\texttt{fixed\_128} & 38.3 & 50.7 & \textbf{95.0} & \textbf{77.1} & 125.3 & 631 \\
\texttt{fixed\_256} & 38.3 & 56.7 & 88.3 & 73.8 & 244.0 & 322 \\
\texttt{fixed\_512} & 40.0 & 53.7 & 90.0 & 63.7 & 472.2 & 164 \\
\texttt{recursive\_256} & \textbf{41.7} & \textbf{60.6} & 91.7 & 72.5 & 175.8 & 389 \\
\texttt{sentence\_256} & 35.0 & 51.0 & 86.7 & 72.1 & 223.9 & 302 \\
\texttt{semantic\_256} & \textbf{41.7} & 56.9 & \textbf{95.0} & 74.6 & 144.8 & 467 \\
\bottomrule
\end{tabular}
\caption{Results on SQuAD v2. Recursive 256 is the strongest end-to-end setting in the current run.}
\label{tab:squad}
\end{table*}

\begin{table*}[t]
\centering
\small
\begin{tabular}{lrrrrrr}
\toprule
System & EM & F1 & Recall@4 & Precision@4 & Avg chunk tokens & \# chunks \\
\midrule
\texttt{parametric\_only} & 10.0 & 27.7 & 0.0 & 0.0 & -- & -- \\
\texttt{fixed\_128} & 20.0 & 36.5 & 75.0 & \textbf{42.5} & 89.3 & 421 \\
\texttt{fixed\_256} & 20.0 & 34.3 & 75.0 & 40.8 & 114.2 & 313 \\
\texttt{fixed\_512} & 20.0 & 35.7 & \textbf{76.7} & 41.7 & 117.4 & 301 \\
\texttt{recursive\_256} & \textbf{26.7} & \textbf{42.6} & \textbf{76.7} & \textbf{42.5} & 113.6 & 313 \\
\texttt{sentence\_256} & 23.3 & 37.0 & \textbf{76.7} & 41.7 & 112.7 & 313 \\
\texttt{semantic\_256} & 23.3 & 42.4 & \textbf{76.7} & \textbf{42.5} & 97.2 & 363 \\
\bottomrule
\end{tabular}
\caption{Results on HotpotQA. Doc-level Recall@4 shows that recovering all supporting documents remains incomplete in the current run.}
\label{tab:hotpot}
\end{table*}

\noindent\textit{Single-run percentages. Current subsets: SQuAD v2 ($n=60$), HotpotQA ($n=30$).}


The parametric-only baseline is weak on both datasets, confirming retrieval sensitivity. On SQuAD, \texttt{recursive\_256} leads at 60.6 F1 and 41.7 EM, with \texttt{semantic\_256} competitive at 56.9 F1 and tied for the best Recall@4 (95.0). \texttt{fixed\_128} achieves the highest retrieval precision (77.1) but lower F1, suggesting that very small chunks recover evidence but may fragment it. Larger fixed windows (\texttt{fixed\_512}) degrade both EM and F1, consistent with the hypothesis that broad chunks dilute relevant evidence with distractors.

On HotpotQA, \texttt{recursive\_256} is best at 42.6 F1 and 26.7 EM, with \texttt{semantic\_256} essentially tied on F1 at 42.4. Doc-level Recall@4 tops out at 76.7 rather than saturating, suggesting that evidence composition---not chunk boundary placement---is the main bottleneck once retrieval is only partially successful.

\section{Discussion}

The results support the proposal's hypothesis that chunking is a meaningful variable, though the conclusions remain preliminary given the small subsets.

\textbf{Chunking matters, and the best strategy depends on question type.} On SQuAD, chunkers that preserve moderately large coherent spans (recursive, semantic) outperform both very small fixed windows and sentence-level segmentation. On HotpotQA, exact match is flatter across systems but F1 still varies, suggesting that multi-hop QA benefits from chunkers preserving local structure. This is consistent with the proposal's expectation that local factoid QA and multi-hop QA reward different chunking behaviors.

\textbf{Generator alignment interacts with chunking quality.} The current SCC rerun with improved prompting and answer compression is substantially stronger than our first Mistral endpoint pass, confirming that generator interface design can move end-to-end QA quality without changing retrieval. Once the prompt is tuned for short grounded spans, recursive chunking becomes the most robust setting across both datasets.

\textbf{Simple baselines are informative.} The parametric-only baseline is far weaker than every retrieval-based system, confirming that the task is genuinely retrieval-sensitive. Among fixed-window baselines, smaller chunks perform better on SQuAD (supporting the claim that large windows introduce distractors), but multi-hop HotpotQA does not consistently reward finer segmentation, because very small chunks can separate evidence that needs to stay together.

\textbf{Failure taxonomy.} A preliminary inspection of predictions reveals three failure categories matching the proposal's taxonomy: retrieval failures (the relevant chunk was never retrieved), context-fragmentation errors (evidence is split unhelpfully across chunks), and generator-interface failures (the correct evidence is present but the model produces a differently worded answer or misses the span entirely).

\subsection{Limitations}

\textbf{Strict metrics and answer-format mismatch.}  EM gives zero credit if the predicted answer contains even one extra word, and F1 gives only partial credit. Mistral tends to produce full-sentence answers rather than short spans, so a portion of the current score gap may reflect formatting mismatches rather than true content errors. A preliminary inspection suggests that a meaningful fraction of EM~=~0 cases represent correct answers worded differently from the gold label (e.g., ``the Denver Broncos'' vs.\ ``Denver Broncos'').

\textbf{Uniform prompting across datasets.}  The system uses a single extractive QA prompt for both datasets. HotpotQA's multi-hop reasoning may benefit from a different prompt strategy and from including document titles in the context so the model can tell which passage each fact comes from.

\textbf{Limited post-processing.}  The existing answer normalizer handles common formatting issues but does not yet cover all observed mismatch patterns such as determiners, parenthetical clarifications, and unit formatting.

\textbf{Small evaluation subsets.}  The subsets (60 and 30 questions) are sufficient to reveal broad trends, but the ranking among chunkers could shift with larger samples.

\subsection{Remaining Work}

The remaining work directly addresses these limitations:
\begin{enumerate}
  \item \textbf{Failure analysis:} categorize all EM~=~0 predictions into retrieval failures, generation errors, and surface-form mismatches to quantify how much of the score gap is real vs.\ formatting.
  \item \textbf{Targeted improvements:} introduce dataset-specific prompts for HotpotQA, add document titles to context, and extend answer normalization based on observed mismatch patterns.
  \item \textbf{Scale experiments:} rerun on larger subsets with multiple seeds to confirm the observed chunking trends are stable.
  \item \textbf{Strengthen diagnostics:} separate retrieval failure from generation failure more cleanly, especially for HotpotQA where multi-hop reasoning makes attribution harder.
  \item \textbf{Broaden retrieval:} test whether chunking effects persist under BM25, hybrid, and reranked retrieval settings already implemented in the codebase.
\end{enumerate}

Relative to the proposal, the project has stayed on the same conceptual trajectory. The main new insight is that switching to a stronger 7B-class model is useful only when the generator interface is tuned for grounded QA; once that mismatch is reduced, the chunking comparison becomes substantially cleaner and more useful for final-stage analysis.

\bibliography{references}
\bibliographystyle{acl_natbib}

\end{document}
